%=============================================================================
%  DeepCardio-RAG — Genuine Training Pipeline: Technical Specification Report
%
%  Companion document to DeepCardio_Genuine_Training.ipynb and to
%  DeepCardio_CompIntel_Submit.tex (the current paper version).
%
%  PROVENANCE RULE FOR THIS DOCUMENT:
%  Every architectural figure is read from the SOURCE CODE, and every
%  performance figure from data/genuine_results.json (+ data/vfdb_cv_metrics.json).
%  Nothing is copied from a paper draft. Where the paper and the code
%  disagree, the code value is given in the body and the disagreement is
%  recorded in Appendix A.
%
%  Compile:  pdflatex DeepCardio_Genuine_Training_TechReport.tex   (x2)
%=============================================================================
\documentclass[11pt,a4paper]{article}

\usepackage[margin=2.2cm]{geometry}
\usepackage[T1]{fontenc}
\usepackage{lmodern}
\usepackage{booktabs}
\usepackage{longtable}
\usepackage{array}
\usepackage{amsmath}
\usepackage{amssymb}
\usepackage{xcolor}
\usepackage{graphicx}
\usepackage{enumitem}
\usepackage{fancyhdr}
\usepackage{titlesec}
\usepackage[hidelinks,bookmarks=true]{hyperref}

\definecolor{dcblue}{RGB}{18,72,120}
\definecolor{dcgrey}{RGB}{95,105,115}
\definecolor{dcwarn}{RGB}{150,60,20}
\definecolor{rowbg}{RGB}{243,246,249}

\titleformat{\section}{\normalfont\Large\bfseries\color{dcblue}}{\thesection}{0.7em}{}
\titleformat{\subsection}{\normalfont\large\bfseries\color{dcblue}}{\thesubsection}{0.7em}{}

\pagestyle{fancy}
\fancyhf{}
\fancyhead[L]{\small\color{dcgrey}DeepCardio-RAG --- Genuine Training Pipeline}
\fancyhead[R]{\small\color{dcgrey}\thepage}
\renewcommand{\headrulewidth}{0.4pt}

\newcommand{\code}[1]{\texttt{\small #1}}
\newcommand{\yes}{\textcolor{dcblue}{\textbf{yes}}}
\newcommand{\no}{\textcolor{dcwarn}{\textbf{no}}}

% Boxed note environment
\usepackage{mdframed}
\newmdenv[linecolor=dcblue,linewidth=0.9pt,backgroundcolor=rowbg,
          innertopmargin=6pt,innerbottommargin=6pt,skipabove=8pt,skipbelow=8pt]{notebox}

% Long \texttt{} identifiers cannot hyphenate; prefer looser interword
% spacing over text running into the margin.
\sloppy
\emergencystretch=3em

\begin{document}

%-----------------------------------------------------------------------------
\begin{titlepage}
\centering
\vspace*{2.2cm}
{\Huge\bfseries\color{dcblue} DeepCardio-RAG\par}
\vspace{0.6cm}
{\LARGE Genuine Training Pipeline\par}
\vspace{0.35cm}
{\large Technology Stack, Datasets, Algorithms,\\[2pt]
and Input\,/\,Process\,/\,Output Specification\par}
\vspace{1.1cm}
{\large Technical Specification Report\par}
\vspace{0.25cm}
{\color{dcgrey}Companion to \code{DeepCardio\_Genuine\_Training.ipynb}\par}
\vspace{2.0cm}

\begin{minipage}{0.86\textwidth}
\begin{notebox}
\textbf{Provenance of the figures in this document.} Architecture and
hyperparameter values are read from the implementation source
(\code{core/}, \code{colab\_train\_genuine.py}); performance values are read
from \code{data/genuine\_results.json} and
\code{data/vfdb\_cv\_metrics.json}, which store held-out evaluations
produced by the notebook itself. No value is copied from a manuscript
draft. Where a paper draft and the code disagree, the \emph{code} value
appears in the body and the disagreement is recorded in
Appendix~\ref{app:delta}.

\medskip
\textbf{Reporting policy.} Only genuinely trained modules on benchmark data
are reported. Modules without trained weights are marked
\textsc{not trained} and carry no performance figures; no demonstration-patient
or synthetic-signal output is presented as a result.
\end{notebox}
\end{minipage}

\vfill
{\color{dcgrey}\small Research prototype --- not cleared for clinical use.\par}
\vspace{0.8cm}
\end{titlepage}

\tableofcontents
\newpage

%=============================================================================
\section{Purpose and Scope}
%=============================================================================

\code{DeepCardio\_Genuine\_Training.ipynb} is the single, current training
notebook for the DeepCardio-RAG system. Its purpose is narrow and
deliberate: train every module that can be trained on a
\emph{public benchmark dataset}, evaluate it on a \emph{held-out} split,
and persist both the weights and the measured metrics to disk. It exists
to replace an earlier situation in which several modules ran on
random-initialised weights while still emitting clinical-looking output.

This report documents four things, in the order requested:

\begin{enumerate}[leftmargin=1.4em,itemsep=2pt]
  \item the complete \textbf{technology stack} (Section~\ref{sec:stack});
  \item the \textbf{datasets} used, with provenance, size, split protocol
        and known integrity issues (Section~\ref{sec:data});
  \item the \textbf{algorithms} implemented per module
        (Section~\ref{sec:algo});
  \item the \textbf{input\,$\rightarrow$\,process\,$\rightarrow$\,expected
        output} specification for each module, including the observed
        held-out values (Section~\ref{sec:io}).
\end{enumerate}

The notebook does \emph{not} serve the web dashboard. The FastAPI
application (\code{main.py}) is a separate local process; the notebook and
the dashboard share only the artefacts written into \code{data/}.

%=============================================================================
\section{Technology Stack}
\label{sec:stack}
%=============================================================================

\subsection{Runtime topology}

The notebook is \emph{dual-runtime}. Cells are tagged for one of two
execution backends, and the two halves never communicate during a run.

\begin{table}[!ht]
\centering
\caption{Execution backends and their responsibilities}
\label{tab:runtime}
\small
\begin{tabular}{@{}p{3.2cm}p{5.3cm}p{6.6cm}@{}}
\toprule
\textbf{Backend} & \textbf{Hardware} & \textbf{Responsibility}\\
\midrule
Colab cloud & NVIDIA T4 GPU, Google Drive mounted at
\code{/content/drive} &
All model training, dataset download, ablation studies. Cannot reach the
private vector database.\\
\addlinespace
Local runtime & Windows host that also runs the Milvus container;
Jupyter exposed over \code{jupyter\_http\_over\_ws} &
Vector-database checks, retrieval grounding, pinned Milvus client.\\
\bottomrule
\end{tabular}
\end{table}

The split works without any network tunnel because \emph{no training
module touches the vector database} --- the training code contains no
reference to \code{db\_manager}, \code{init\_db} or \code{pymilvus}. The
hand-off between the two halves is Google Drive itself: weights written to
\code{data/} on Colab synchronise back down to the local project folder.
The largest artefact, \code{ecg\_ptbxl\_encoder.pt}, is approximately
55~MB, so the local half must wait for that synchronisation to finish
before it can confirm that the pipeline is using trained weights.

\subsection{Software stack}

\small
\begin{longtable}{@{}p{2.9cm}p{4.1cm}p{8.3cm}@{}}
\caption{Software components and their role in the training pipeline}
\label{tab:stack}\\
\toprule
\textbf{Layer} & \textbf{Component} & \textbf{Role}\\
\midrule
\endfirsthead
\multicolumn{3}{@{}l}{\small\itshape Table~\ref{tab:stack} continued}\\
\toprule
\textbf{Layer} & \textbf{Component} & \textbf{Role}\\
\midrule
\endhead
\bottomrule
\endfoot
Deep learning & PyTorch (\code{torch}, \code{torch.nn}, \code{torch.optim})
& All eight neural modules; CUDA on Colab, CPU fallback locally.\\
& \code{torch.utils.data} & \code{DataLoader} / \code{TensorDataset} /
custom \code{Dataset} for the video and audio paths.\\
\addlinespace
Language models & \code{transformers==4.44.2} & GPT-2 (124M) report decoder
with soft-prompt injection; pinned because a newer build broke
\code{GPT2LMHeadModel} import on Colab.\\
& \code{sentence-transformers} & Sentence-BERT \code{all-mpnet-base-v2},
768-dim embeddings for retrieval.\\
\addlinespace
Classical ML & \code{scikit-learn} &
\code{RepeatedStratifiedKFold}, \code{StandardScaler}, \code{KNNImputer},
\code{GradientBoostingClassifier}, \code{RandomForestClassifier},
\code{ExtraTreesClassifier}, \code{LogisticRegression}, and all metrics
(accuracy, macro-$F_1$, ROC-AUC, Brier).\\
& \code{joblib} & Persists fitted scalers and feature-column lists
alongside the network weights.\\
\addlinespace
Signal / data I/O & \code{wfdb>=4.3.1} & Decodes PhysioNet WFDB waveforms
and annotations (PTB-XL \code{.hea}/\code{.dat}, VFDB \code{.atr}).
\textbf{Version-critical} --- see Section~\ref{sec:pins}.\\
& \code{pandas==2.2.3} & Tabular loaders (NHANES, UCI Cleveland, MIT-BIH
CSV). Pinned by the \code{wfdb} requirement.\\
& \code{numpy} (2.x) & Array backbone throughout.\\
& \code{opencv-python-headless} & Echo video decoding and frame resizing.\\
& \code{kagglehub} & Programmatic Kaggle dataset download (PTB-XL, NHANES).\\
\addlinespace
Vector database & \code{pymilvus==2.6.9} & Client for Milvus standalone
\code{v2.6.2} at \code{localhost:19530}. \textbf{Version-critical}.\\
& Milvus standalone (Docker) & Collection
\code{cardio\_knowledge\_base}: 768-dim, $L_2$ metric, \code{IVF\_FLAT}
index, \code{nlist}~=~2048.\\
\addlinespace
Serving (separate) & \code{fastapi} + \code{uvicorn} & REST API and static
dashboard; not exercised by the notebook.\\
& \code{fpdf2} & Dual-audience PDF report generation.\\
\addlinespace
Configuration & \code{pydantic-settings} & \code{config.py} settings object;
real environment variables take precedence over \code{.env}.\\
\end{longtable}

\subsection{Version pins that are load-bearing}
\label{sec:pins}

Three pins in the dependency cell are not cosmetic. Each was added in
response to a specific, diagnosed failure.

\begin{table}[!ht]
\centering
\caption{Load-bearing version pins}
\label{tab:pins}
\small
\begin{tabular}{@{}p{3.0cm}p{11.8cm}@{}}
\toprule
\textbf{Pin} & \textbf{Failure it prevents}\\
\midrule
\code{wfdb>=4.3.1} &
Under NumPy~2, older \code{wfdb.rdann()} raises
\code{OverflowError: Python integer 256 out of bounds for uint8}. The VFDB
loader then fell back to a hand-rolled byte parser whose SKIP handling
byte-swapped the 32-bit interval, yielding annotation offsets up to
$1.99\times10^{10}$ on a 525{,}000-sample signal. Every window fell before
the first valid annotation and was labelled Normal, producing a
\emph{single-class} dataset that scored accuracy 1.0 --- a trivial
solution, not a result.\\
\addlinespace
\code{pandas==2.2.3} &
Required by \code{wfdb} 4.3.1. Colab's baseline is 2.2.2; a blanket
\code{--upgrade} pulls pandas 3.0.x and breaks the pandas-heavy tabular
pipelines. The resulting
\code{google-colab 1.0.0 requires pandas==2.2.2} warning is expected and
harmless --- it governs Colab's own table display, not training.\\
\addlinespace
\code{pymilvus==2.6.9} &
\code{pymilvus} 3.x removes the 2.x ORM API
(\code{connections} / \code{Collection} / \code{FieldSchema}) that
\code{database/db\_manager.py} is built on, and pulls
\code{setuptools}~83, violating \code{torch}'s \code{setuptools<82}
requirement.\\
\addlinespace
\code{torch} (not upgraded) &
Colab's preinstalled CUDA build works; upgrading it breaks
\code{transformers} model imports (e.g.\
\code{Could not import GPT2LMHeadModel}).\\
\bottomrule
\end{tabular}
\end{table}

\subsection{Notebook execution map}

\begin{table}[!ht]
\centering
\caption{Cell-to-runtime map (0-indexed over all cells, markdown included)}
\label{tab:cells}
\small
\begin{tabular}{@{}p{1.6cm}p{2.4cm}p{10.6cm}@{}}
\toprule
\textbf{Cell} & \textbf{Runtime} & \textbf{Function}\\
\midrule
2  & either & Runtime detection; mounts Drive on Colab, sets
\code{MILVUS\_HOST}/\code{MILVUS\_PORT} locally; inserts CWD on
\code{sys.path}.\\
4  & Colab & Dependency install with the pins of Table~\ref{tab:pins}.\\
5  & either & In-kernel version assertion (pip output alone is not proof
the kernel took the upgrade).\\
7  & Colab & Kaggle credentials.\\
9  & Colab & PTB-XL completeness fix; forces a full download to local disk
and repoints the loader.\\
11 & Colab & PhysioNet fetch of VFDB and CirCor via
\code{wget -r -N -c}; completeness check against \code{RECORDS}.\\
13 & Colab & Optional SHA256 byte-level integrity check.\\
15 & Colab & \code{run\_all()} --- trains every module in sequence.\\
17--19, 21 & Colab & Individual modules: PTB-XL flagship, arthritis,
MIT-BIH arrhythmia, CardioFusion.\\
20 & local & Pins \code{pymilvus} to 2.6.9.\\
23--26 & Colab & Echo, VFDB, PCG, heart disease.\\
28--29 & local & Milvus connection and schema check; confirms the project
is on the real server and not a silent fallback.\\
31 & local & Confirms the pipeline loads the trained encoder and that
retrieval is diagnosis-grounded.\\
33 & either & RAG retrieval evaluation (in-memory).\\
35 & Colab & Ablation studies.\\
37 & Colab & VFDB recording-level cross-validation.\\
\bottomrule
\end{tabular}
\end{table}

\begin{notebox}
\textbf{Two operational cautions carried in the notebook itself.}
(i)~Cell~4 pins \emph{Colab's} environment; it should not be re-run on the
local runtime, where it would churn a working install. On a runtime
switch, re-run cell~2 only.
(ii)~Section~5c must be skipped on a Colab cloud backend: a cloud VM
cannot route to a private machine, and \code{db\_manager} will silently
fall back to FAISS while still logging \code{Vector DB ready}. That green
message is not evidence of anything.
\end{notebox}

%=============================================================================
\section{Datasets}
\label{sec:data}
%=============================================================================

Seven public corpora are used. Table~\ref{tab:datasets} gives the summary;
the subsections record provenance, split protocol and the integrity
issues that were found and fixed.

\begin{table}[!ht]
\centering
\caption{Datasets used by the genuine training pipeline}
\label{tab:datasets}
\small
\begin{tabular}{@{}p{2.7cm}p{3.4cm}p{2.9cm}p{2.0cm}p{3.0cm}@{}}
\toprule
\textbf{Module} & \textbf{Dataset} & \textbf{Size used} &
\textbf{Classes} & \textbf{Split protocol}\\
\midrule
Flagship ECG & PTB-XL 500\,Hz (PhysioNet) & 3\,000 records sampled
(2\,389/296/315) & 5 superclasses & Official folds 1--8/9/10,
patient-independent\\
\addlinespace
ECG arrhythmia & MIT-BIH AAMI beats (Kaggle CSV) & 2\,999 train /
21\,892 test & 5 (N,S,V,F,Q) & Official train/test CSV split\\
\addlinespace
CardioFusion & MIT-BIH, ECG-signal path only & 2\,549 train /
21\,892 test & 5 & Train/val carve-out; test scored once\\
\addlinespace
Ventricular & MIT-BIH VFDB (PhysioNet) & 22 recordings,
5\,s windows & 2 binary + 4 rhythm & Whole recordings per split;
$2\times5$-fold recording-level CV\\
\addlinespace
Heart sound & CirCor DigiScope (PhysioNet) & 942 subjects,
2\,057/480/626 recordings & 3 murmur & Subject-level
(patient-independent)\\
\addlinespace
Heart disease & UCI Cleveland, Kaggle mirror & 1\,025 rows
$\rightarrow$ \textbf{302 unique} & 2 & $5\times3$ repeated stratified
K-fold + 25\% hold-out\\
\addlinespace
Arthritis & NHANES MCQ (CDC/Kaggle) & 5\,756 records
(3\,000 sampled) & 2 & $5\times3$ repeated stratified K-fold + 20\%
hold-out\\
\addlinespace
Echo (LVEF) & EchoNet-Dynamic (Stanford AIMI) &
\textcolor{dcwarn}{not obtained} & 3 EF categories &
\textcolor{dcwarn}{n/a --- access required}\\
\bottomrule
\end{tabular}
\end{table}

\subsection{PTB-XL --- flagship 12-lead ECG}

Clinical 12-lead ECGs at 500~Hz, 10~seconds per record
($12\times5000$ samples). The loader targets the five diagnostic
superclasses \{NORM, MI, STTC, CD, HYP\} and uses the official
patient-independent folds, so no patient appears in more than one split.

PTB-XL assigns \emph{multi-label} superclasses, so the published
per-class record counts overlap rather than partition the corpus:
NORM 9\,528, MI 5\,486, STTC 5\,250, CD 4\,907, HYP 2\,655 over 21\,837
records from 18\,885 patients. These are reference frequencies stored in
the loader as \code{PTBXL\_REFERENCE\_COUNTS}; the training run samples a
subset (3\,000 records by default) and reports metrics on its own held-out
fold.

\emph{Integrity note.} Drive-cached copies of PTB-XL are sometimes missing
waveform folders, and a missing record silently falls back to a
\emph{synthetic} signal --- which would contaminate results with fabricated
data. Cell~9 exists solely to prevent this: it force-downloads a complete
copy to local Colab disk and repoints the loader. The acceptance check is
that \code{records500/21000} contains approximately 838 \code{.hea} files;
a count of~0 means the mirror itself is incomplete.

\subsection{MIT-BIH Arrhythmia --- beat classification}

Pre-segmented single-beat vectors of 187 samples, distributed as CSV, with
the five AAMI classes N, S, V, F, Q. The official train/test CSV split is
used; the test set of 21\,892 beats is many times larger than the 2\,999
sampled training beats, which makes the reported accuracy pessimistic
relative to a balanced protocol. The class distribution is strongly
imbalanced toward N, so macro-$F_1$ and AUC are the honest headline
figures and raw accuracy is not.

\subsection{MIT-BIH VFDB --- ventricular arrhythmia}

Twenty-two dual-channel recordings, 30 minutes each at 250~Hz, segmented
into 5-second windows. Rhythms \{VT, VF, VFL, VFIB, ASYS, HGEA, VER\} are
treated as \emph{dangerous}; the remainder (N, AFIB, NOISE, SBR, SVTA,
NOD, B, PM) form the background class. Splits are by whole recording, so a
recording never contributes windows to two splits.

\emph{Integrity note.} Twenty-two recordings cannot support a point
estimate from a single 15/3/4 split. The single split gave test
$F_1$(dangerous)~0.5147 / AUC~0.8046 against validation 0.894 / 0.984 ---
a gap over four test recordings. It was therefore replaced by
$2\times5$-fold recording-level cross-validation (cell~37), and the
cross-validated figures are the reportable ones.

\subsection{CirCor DigiScope --- phonocardiogram murmur}

Heart-sound recordings at 4\,000~Hz from multiple auscultation locations
(AV, PV, TV, MV, Phc), labelled \{Absent, Present, Unknown\}. \emph{Unknown}
is a genuine annotation meaning auscultation was inconclusive --- not a
missing-data marker. The public training subset is 3\,163 recordings from
942 subjects; the full release (5\,272 recordings, 1\,568 subjects) retains
a hidden test set.

\emph{Integrity notes, both fixed.} (i)~Subject identifiers were parsed
with \code{rsplit("\_", 1)}, which splits repeat-recording names such as
\code{50204\_MV\_1.wav} on the \emph{last} underscore, yielding patient id
\code{50204\_MV} instead of \code{50204}. This created 20 phantom subjects
(962 counted against 942 real), silently defaulted their murmur label to
the \emph{Unknown class}, and --- most seriously --- allowed one patient to
appear in both train and test. Corrected to \code{partition("\_")}.
(ii)~A recording truncated by an interrupted download
(\code{85286\_AV.wav}, 15\,974 of 161\,280 bytes) was repaired; the
download command uses \code{wget -N -c} and never \code{-nc}, which skips
any file that merely exists without checking completeness.

\subsection{UCI Cleveland --- heart disease}

The Kaggle mirror ships 1\,025 rows, but these are the 303-record Cleveland
cohort expanded with duplicate rows. Identical rows on both sides of a
random split inflate every metric, so training deduplicates to the
\textbf{302 unique} records first.

\emph{Integrity note.} This mirror's \code{target} column is inverted
relative to the module's own documentation: measured on the file, the
\code{target==0} group carries the diseased clinical profile. The
prediction path detects the polarity from clinical priors rather than
hardcoding a flip, so a corrected upstream file would be left alone.

\subsection{NHANES --- arthritis risk}

CDC National Health and Nutrition Examination Survey, Medical Conditions
Questionnaire: 5\,756 records (2\,752 male, 3\,004 female, mean age
49.1~years). The prediction target is \textbf{MCQ160A}, ``has a doctor ever
told you that you have arthritis'' --- self-reported, doctor-diagnosed,
predominantly osteoarthritis. It is \emph{not} serological rheumatoid
arthritis, and results are not comparable with RA-cohort literature.

\emph{Integrity note.} The arthritis-type field MCQ195 (RA/OA) is a
skip-pattern question asked only of arthritis-positive respondents; it is
excluded as a target-leakage source. Retaining it inflated AUC from an
honest $\approx$0.79 to $\approx$0.99 --- the ablation study
that revealed this is cell~35.

\subsection{EchoNet-Dynamic --- not obtained}

Apical 4-chamber echocardiogram videos with EF, ESV and EDV annotations.
Access requires registration with Stanford AIMI and the dataset is
\textbf{not present}; \code{genuine\_results.json} records an \code{error}
entry for this module rather than a metric. The echo module therefore has
\textbf{no trained weights}, and every figure that would describe echo
performance is withheld throughout this report.

%=============================================================================
\section{Algorithms}
\label{sec:algo}
%=============================================================================

All architectural values below are read from the implementation.

\subsection{Flagship encoder --- \texorpdfstring{\code{ECGTransformerMoE}}{ECGTransformerMoE}}

Defined in \code{core/pipeline.py}; trained by \code{core/train\_ptbxl.py}
through the wrapper \code{ECGSuperclassNet}.

\begin{enumerate}[leftmargin=1.5em,itemsep=1pt]
  \item \textbf{Patch embedding}: \code{Conv1d(12, 256, kernel=50,
        stride=50)} maps the $12\times5000$ input to 100 non-overlapping
        patches of 256 dimensions --- one patch per 100~ms at 500~Hz.
  \item A learnable \code{[CLS]} token is prepended, giving
        $(B, 101, 256)$, and learnable positional embeddings are added.
  \item \textbf{6 Transformer blocks}, each with 8-head self-attention and
        a \textbf{mixture-of-experts feed-forward} layer: 4 experts,
        top-2 routing, FFN dimension 1024, dropout 0.1.
  \item The \code{[CLS}] output passes through LayerNorm and a
        $256\rightarrow256$ projection to yield the ECG embedding.
  \item A diagnostic head maps the embedding to the 5 PTB-XL superclasses.
        The embedding is separately projected $256\rightarrow768$ to query
        the vector store.
\end{enumerate}

Total parameters: 14\,453\,765. Optimiser Adam, learning rate
$3\times10^{-4}$, batch size 32, 30 epochs, seed 42. A ResNet-34 1D
encoder is available as the ablation baseline
(\code{ablate\_ecg\_encoder}).

\subsection{Beat classifier --- \texorpdfstring{\code{ECGArrhythmiaCNN}}{ECGArrhythmiaCNN}}

Four 1D-convolution blocks with channel widths $(32, 64, 128, 256)$ and
kernels $(7,5,3,3)$, each followed by batch normalisation, ReLU and
pooling; then a fully connected head $256\rightarrow128\rightarrow64
\rightarrow5$. The input tensor is single-channel,
$(N, 1, 187)$. Class imbalance is handled with a weighted
cross-entropy loss.

\subsection{Ventricular detector --- \texorpdfstring{\code{VentricularArrhythmiaDetector}}{VentricularArrhythmiaDetector}}

A dual-channel 1D-CNN (\code{in\_channels=2}) over 5-second windows, with
\emph{two heads}: a binary dangerous-vs-background head and a 4-class
rhythm-family head \{Normal, VT, VF/VFL, Other\_Dangerous\}. 313\,958
parameters. The class carries an explicit \code{is\_trained} flag, set
true only when genuine weights are loaded, so that a random-weight
detector cannot surface a live alert.

Cross-validation (\code{core/cv\_vfdb.py}) folds \emph{whole recordings},
stratified by dangerous-window fraction; validation recordings are drawn
from the training pool only, and each test fold is scored exactly once
after early stopping. Single-class test folds report AUC as NaN and are
counted rather than silently averaged; a single-class \emph{training} fold
raises.

\subsection{Murmur classifier --- \texorpdfstring{\code{HeartSoundClassifier}}{HeartSoundClassifier}}

Audio at 4\,000~Hz is cut into 5-second segments and converted to log-Mel
spectrograms with 64 Mel bands and hop length 128, padded or cropped to a
fixed width of 160 frames for batching. A 2D-CNN with three convolution
blocks $(16, 32, 64)$ and adaptive average pooling feeds a 3-class softmax
head. 44\,707 parameters --- by far the smallest network in the system.

\subsection{Tabular BERT + Mixture-of-Experts}

Two instances of the same design, with \emph{different} hyperparameters:

\begin{table}[!ht]
\centering
\caption{TabularBERT-MoE instances as implemented}
\label{tab:bertmoe}
\small
\begin{tabular}{@{}p{4.2cm}p{2.2cm}p{2.0cm}p{2.0cm}p{2.2cm}@{}}
\toprule
\textbf{Instance} & \textbf{$d_{\text{model}}$} & \textbf{Heads} &
\textbf{Layers} & \textbf{Experts}\\
\midrule
Arthritis (\code{TabularBERTMoE}) & 64 & 4 & 3 & 6\\
Heart disease (\code{HeartDiseaseBERTMoE}) & 64 & 4 & 2 & 4\\
\bottomrule
\end{tabular}
\end{table}

Each scalar feature value is projected to $d_{\text{model}}$ and summed
with a learned feature-identity embedding, so the feature set is treated
as a token sequence; a Transformer encoder then contextualises the tokens
and a gated mixture-of-experts head produces the class logits.

For arthritis the network is one input to a \textbf{stacking ensemble}:
Gradient Boosting, Random Forest and Extra Trees form the base layer
alongside the BERT-MoE, and a logistic-regression meta-learner is fitted
on out-of-fold predictions. Minority-class balancing (SMOTE-style
oversampling), imputation and scaling are all fitted \emph{inside} each
training fold, never across the split.

For heart disease, plain logistic-regression and gradient-boosting
baselines are evaluated on \emph{identical folds}, so the added complexity
of the deep model can be judged against them.

\subsection{CardioFusion}

A cross-modal Transformer with a multi-gate mixture-of-experts head over
six tasks. Multi-task loss weights are \emph{learned} by homoscedastic
uncertainty weighting (Kendall et al.), not hand-set. Because no public
benchmark provides same-patient four-modality paired data, only the
\textbf{ECG-signal path} is trained, on MIT-BIH beats resampled to a
512-sample segment length. The remaining heads and the joint contrastive
alignment are untrained --- visible in the learned task weights, where
only the arrhythmia task has moved off its 0.5 initialisation
(to 0.4196).

\subsection{Retrieval subsystem}

Retrieval is \emph{diagnosis-grounded}: the query is built from the
predicted PTB-XL superclass, not from a raw embedding, and if the
diagnostic head is absent the pipeline deliberately serves labelled
placeholders rather than querying with something unfounded.

Three rankers are evaluated over a 67-document clinical corpus with 38
labelled queries: a lexical TF-IDF baseline, a semantic Sentence-BERT
ranker (\code{all-mpnet-base-v2}, 768-dim), and a \textbf{hybrid
reciprocal-rank-fusion} ranker combining the two. Metrics are
Recall@$k$, Precision@$k$, MRR, nDCG@$k$ and Hit@$k$, reported overall and
per superclass.

%=============================================================================
\section{Input, Process and Expected Output}
\label{sec:io}
%=============================================================================

This section is the operational specification: for each module, what goes
in, what the algorithm does, and what should come out --- both in
\emph{form} (the shape and keys of the result) and in \emph{value} (the
held-out figures actually observed).

\begin{notebox}
\textbf{How to read the ``expected value'' column.} These are the
measurements stored in \code{data/genuine\_results.json} from the runs
that produced the current weights. A re-run on the same seed and record
count should land close to them; a substantially \emph{higher} figure is a
reason to check for leakage, not a success.
\end{notebox}

\subsection{Module 1 --- Flagship ECG superclass encoder}

\begin{table}[!ht]
\centering
\small
\begin{tabular}{@{}p{2.4cm}p{12.4cm}@{}}
\toprule
\textbf{Entry point} & \code{train\_ptbxl(n\_records=3000, epochs=30)}
--- notebook cell~17\\
\midrule
\textbf{Input} & PTB-XL WFDB records; tensor $(B, 12, 5000)$ --- 12 leads,
10~s at 500~Hz. Labels: 5 diagnostic superclasses. Patient-independent
official folds 1--8 (train), 9 (val), 10 (test).\\
\addlinespace
\textbf{Process} & Patch embedding (100 patches of 50 samples)
$\rightarrow$ \code{[CLS]} + positional embeddings $\rightarrow$ 6
Transformer blocks with 4-expert top-2 MoE FFN $\rightarrow$ LayerNorm
$\rightarrow$ 256-dim embedding $\rightarrow$ diagnostic head. Adam,
lr~$3\times10^{-4}$, batch~32, 30 epochs; model selection on validation
macro-$F_1$; test scored once.\\
\addlinespace
\textbf{Output artefacts} & \code{data/ecg\_ptbxl\_encoder.pt} ($\approx$55~MB,
loaded by \code{core/pipeline.py}), \code{data/ecg\_ptbxl\_metrics.json}\\
\addlinespace
\textbf{Expected values} &
accuracy \textbf{0.5778} $\cdot$ macro-$F_1$ \textbf{0.4866} $\cdot$
macro-AUC (OvR) \textbf{0.8007} $\cdot$ Brier 0.5966 $\cdot$ ECE 0.117.
Per-class AUC: NORM 0.8832, MI 0.7817, STTC 0.8380, CD 0.7432, HYP 0.7572.
Best validation macro-$F_1$ 0.5418. 14.45M parameters, 209.4~s on CUDA.\\
\bottomrule
\end{tabular}
\end{table}

\subsection{Module 2 --- MIT-BIH beat arrhythmia classifier}

\begin{table}[!ht]
\centering
\small
\begin{tabular}{@{}p{2.4cm}p{12.4cm}@{}}
\toprule
\textbf{Entry point} & \code{train\_ecg\_arrhythmia(n\_records=3000,
epochs=25)} --- notebook cell~19\\
\midrule
\textbf{Input} & MIT-BIH beat CSV; tensor $(N, 1, 187)$. Labels: 5 AAMI
classes (N, S, V, F, Q). Official train/test split.\\
\addlinespace
\textbf{Process} & Four 1D-conv blocks $(32,64,128,256)$ with batch norm
and pooling $\rightarrow$ FC $256\rightarrow128\rightarrow64\rightarrow5$.
Class-weighted cross-entropy for the heavy N-class imbalance.\\
\addlinespace
\textbf{Output artefacts} & \code{data/ecg\_arrhythmia\_cnn\_model.pt},
\code{data/ecg\_arrhythmia\_scaler.pkl}\\
\addlinespace
\textbf{Expected values} & accuracy 0.6065 $\cdot$ macro-$F_1$ 0.5086
$\cdot$ \textbf{macro-AUC (OvR) 0.9283}, on 21\,892 held-out beats.
Quote AUC and macro-$F_1$; accuracy is depressed by the 2\,999/21\,892
train-to-test ratio.\\
\bottomrule
\end{tabular}
\end{table}

\subsection{Module 3 --- Ventricular arrhythmia detector (VFDB)}

\begin{table}[!ht]
\centering
\small
\begin{tabular}{@{}p{2.4cm}p{12.4cm}@{}}
\toprule
\textbf{Entry point} & \code{train\_vfdb(epochs=30)} --- cell~24;
\code{cross\_validate\_vfdb(n\_folds=5, n\_repeats=2, epochs=30)} ---
cell~37\\
\midrule
\textbf{Input} & 22 VFDB recordings, dual-channel, 250~Hz, cut into
5-second windows $(B, 2, 1250)$. Labels: binary dangerous-vs-background
plus 4-class rhythm family.\\
\addlinespace
\textbf{Process} & 1D-CNN encoder with two heads. Folds are whole
recordings stratified by dangerous-window fraction; validation drawn from
the training pool; each test fold scored once after early stopping on
validation $F_1$(dangerous). Training raises if any split is
single-class.\\
\addlinespace
\textbf{Output artefacts} & \code{data/vfdb\_detector.pt},
\code{data/vfdb\_metrics.json}, \code{data/vfdb\_cv\_metrics.json}\\
\addlinespace
\textbf{Expected values} &
\textbf{Cross-validated (reportable):} AUC \textbf{0.9187}
[0.8466, 0.9907] $\cdot$ $F_1$(dangerous) 0.7184 [0.6217, 0.8151] $\cdot$
accuracy 0.8672 [0.8127, 0.9216] $\cdot$ rhythm macro-$F_1$ 0.4512
[0.3831, 0.5193], over 10 fits, 176.8~s, 0 single-class test folds.
Observed fold range AUC [0.6414, 0.9833].
\textbf{Single split (provenance only):} AUC 0.8046, $F_1$(dang.) 0.5147.\\
\bottomrule
\end{tabular}
\end{table}

\begin{notebox}
\textbf{Interval caveat, to be quoted with the number.} The 95\% figure is
a $t$-interval over fold scores. Folds share training recordings, so the
scores are not independent and the interval is \emph{optimistic}: it
describes spread across folds, not a confidence interval over patients.
The superseded single split scored 0.8046, which lies \emph{below} the
interval low of 0.8466 --- direct evidence that the interval is too
narrow. The rhythm head (0.4512) must not be quoted as a strength; the
binary head is the safety-relevant output.
\end{notebox}

\subsection{Module 4 --- Heart sound murmur (CirCor)}

\begin{table}[!ht]
\centering
\small
\begin{tabular}{@{}p{2.4cm}p{12.4cm}@{}}
\toprule
\textbf{Entry point} & \code{train\_circor(n\_max=3000, epochs=25)} ---
notebook cell~25\\
\midrule
\textbf{Input} & CirCor \code{.wav} at 4\,000~Hz $\rightarrow$ 5-second
segments $\rightarrow$ log-Mel spectrogram $(B, 1, 64, 160)$. Labels:
\{Absent, Present, Unknown\}. Subject-level split,
2\,057/480/626 recordings over 942 subjects.\\
\addlinespace
\textbf{Process} & 2D-CNN $(16,32,64)$ with batch norm and adaptive
average pooling $\rightarrow$ 3-class softmax. Early stopping on
validation macro-$F_1$; test scored once.\\
\addlinespace
\textbf{Output artefacts} & \code{data/heart\_sound\_classifier.pt},
\code{data/circor\_metrics.json}\\
\addlinespace
\textbf{Expected values} & accuracy 0.7604 $\cdot$ macro-$F_1$ 0.5347
$\cdot$ \textbf{Present-vs-Absent AUC 0.7611} $\cdot$ per-class $F_1$:
Absent 0.8616, Present 0.4776, Unknown 0.2651. 44\,707 parameters,
1\,249.3~s.\\
\bottomrule
\end{tabular}
\end{table}

\noindent
\textbf{Acceptance checks for a re-run} (all three passed on the current
weights): the loader must report \textbf{942 subjects}, not 962 --- 962
means a stale loader is cached and the kernel needs restarting; \emph{no}
recordings skipped for a missing metadata row; and Present-vs-Absent AUC
\emph{below} 0.7981, because 0.7981 is the leakage-contaminated figure and
leakage inflates. Validation macro-$F_1$ is unstable across epochs, so
quote the AUC as the headline and do not present macro-$F_1$ as precise.

\subsection{Module 5 --- Heart disease risk (UCI Cleveland)}

\begin{table}[!ht]
\centering
\small
\begin{tabular}{@{}p{2.4cm}p{12.4cm}@{}}
\toprule
\textbf{Entry point} & \code{train\_heartdisease(n\_repeats=5, n\_splits=3,
epochs=80)} --- notebook cell~26\\
\midrule
\textbf{Input} & \code{data/heart.csv} deduplicated 1\,025
$\rightarrow$ 302 unique rows; 13 clinical features expanded to 21 after
encoding. Binary target.\\
\addlinespace
\textbf{Process} & $5\times3$ repeated stratified K-fold with KNN
imputation and scaling fitted inside each fold. TabularBERT-MoE
(2 layers, 4 experts) trained per fold, plus logistic-regression and
gradient-boosting baselines on identical folds. 25\% hold-out evaluated
once.\\
\addlinespace
\textbf{Output artefacts} & \code{data/heart\_disease\_bert\_moe\_model.pt},
\newline\code{data/heart\_disease\_bert\_moe\_scaler.pkl},
\newline\code{data/heart\_disease\_metrics.json}\\
\addlinespace
\textbf{Expected values} &
BERT-MoE: accuracy 0.7967 [0.7797, 0.8136], AUC 0.8718 [0.8594, 0.8841],
$F_1$ 0.7939, Brier 0.1749.
\textbf{Logistic baseline: AUC 0.8921 [0.8861, 0.8980]} --- higher than the
deep model. GBT: AUC 0.8833. Hold-out (25\%) BERT-MoE: accuracy 0.7763,
AUC 0.8551.\\
\bottomrule
\end{tabular}
\end{table}

\noindent
The baseline comparison is the point of this module: on 302 unique records
a plain logistic regression beats the TabularBERT-MoE on every metric
including calibration. That result should be reported as it stands.

\subsection{Module 6 --- Arthritis risk (NHANES)}

\begin{table}[!ht]
\centering
\small
\begin{tabular}{@{}p{2.4cm}p{12.4cm}@{}}
\toprule
\textbf{Entry point} & \code{train\_arthritis(n\_records=3000)} ---
notebook cell~18\\
\midrule
\textbf{Input} & NHANES MCQ; 10 base predictors (Age, Gender\_M,
PovertyRatio, EducationLevel, BMI, SystolicBP, DiastolicBP, HasGout,
HasOsteoporosis, InflammationProxy) expanded to \textbf{17} with
engineered terms (age\_risk, age\_squared, bmi\_squared, obesity\_flag,
age\_gender, pulse\_pressure, map). Target MCQ160A.\\
\addlinespace
\textbf{Process} & $5\times3$ repeated stratified K-fold with imputation,
scaling and SMOTE-style balancing fitted \emph{inside} folds. Base layer
GBM + RF + ExtraTrees + TabularBERT-MoE (3 layers, 6 experts);
logistic-regression meta-learner on out-of-fold predictions. 20\%
stratified hold-out taken before balancing.\\
\addlinespace
\textbf{Output artefacts} & \code{data/arthritis\_bert\_moe\_model.pt},
\code{data/arthritis\_bert\_moe\_scaler.pkl} (carries the
\code{feature\_cols} list)\\
\addlinespace
\textbf{Expected values} &
CV accuracy 0.7622 [0.7540, 0.7704] $\cdot$ \textbf{CV AUC 0.7925
[0.7813, 0.8038]} $\cdot$ CV macro-$F_1$ 0.6844 $\cdot$ CV Brier 0.1624.
Hold-out: accuracy 0.7700, AUC 0.8127, $F_1$ 0.7024. Top features:
age\_risk 0.157, age\_squared 0.128, Age 0.115.\\
\bottomrule
\end{tabular}
\end{table}

\noindent
\textbf{Two caveats to carry with these numbers.} First, the three leading
features are collinear encodings of age, together roughly 40\% of total
importance --- this is substantially an age model. Second, the hold-out AUC
(0.8127) lies \emph{outside} the cross-validated interval
[0.7813, 0.8038]; the cross-validated figure is the headline. A run
returning AUC near 0.99 indicates MCQ195 has re-entered the feature set.

\subsection{Module 7 --- CardioFusion (ECG path)}

\begin{table}[!ht]
\centering
\small
\begin{tabular}{@{}p{2.4cm}p{12.4cm}@{}}
\toprule
\textbf{Entry point} & \code{train\_cardiofusion(n\_records=3000,
epochs=20)} --- notebook cell~21\\
\midrule
\textbf{Input} & MIT-BIH beats reshaped to 512-sample ECG segments;
2\,549 train / 21\,892 test. Only the ECG-signal modality is populated.\\
\addlinespace
\textbf{Process} & Cross-modal Transformer with multi-gate MoE head over
six tasks; loss weights learned by homoscedastic uncertainty weighting.
Model selection on a validation split carved from train; class-weighted
arrhythmia loss; test scored once.\\
\addlinespace
\textbf{Output artefacts} & \code{data/cardiofusion\_weights.pt},
\code{data/cardiofusion\_metrics.json}\\
\addlinespace
\textbf{Expected values} & accuracy 0.8944 $\cdot$ macro-$F_1$ 0.5401
$\cdot$ AUC 0.8384. Learned task weights: arrhythmia \textbf{0.4196}, all
five other tasks still at their 0.5 initialisation.\\
\bottomrule
\end{tabular}
\end{table}

\noindent
Those five unchanged weights are the diagnostic fact: the corresponding
heads never received gradient. Only the arrhythmia output of this module
is a trained result; the EF, murmur, ventricular-alert and rhythm heads
must not be surfaced, and no composite score aggregating them is
meaningful. The run is also unstable across epochs, so a single run is not
a defensible point estimate.

\subsection{Module 8 --- LVEF estimation (blocked)}

\begin{table}[!ht]
\centering
\small
\begin{tabular}{@{}p{2.4cm}p{12.4cm}@{}}
\toprule
\textbf{Entry point} & \code{train\_echonet(n\_train=3000, n\_eval=1000,
epochs=25)} --- notebook cell~23\\
\midrule
\textbf{Input} & EchoNet-Dynamic A4C videos, 32 frames at
$112\times112$ grayscale. \textcolor{dcwarn}{Dataset not present.}\\
\addlinespace
\textbf{Process} & 3D-CNN video encoder (\code{EchoNetModel},
384-dim embedding) $\rightarrow$ EF regression head, with predictions
clamped to $[0,100]\%$, plus a 3-way EF category head
(HFrEF $<40$, HFmrEF 40--50, Normal $\geq50$).\\
\addlinespace
\textbf{Expected output} & \textcolor{dcwarn}{\textbf{None.}} The module
records an \code{error} entry in \code{genuine\_results.json} and produces
no weights. Consequently the serving layer returns
\code{ef\_predicted: null} with \code{model\_trained: false} rather than a
number.\\
\bottomrule
\end{tabular}
\end{table}

\subsection{Retrieval evaluation}

\begin{table}[!ht]
\centering
\small
\begin{tabular}{@{}p{2.4cm}p{12.4cm}@{}}
\toprule
\textbf{Entry point} & \code{run\_rag\_evaluation(include\_semantic=True)}
--- notebook cell~33\\
\midrule
\textbf{Input} & 67-document clinical guideline corpus; 38 labelled
queries with per-superclass relevance annotations.\\
\addlinespace
\textbf{Process} & Three rankers scored on identical queries: lexical
TF-IDF, semantic Sentence-BERT (\code{all-mpnet-base-v2}), and hybrid
reciprocal rank fusion. Metrics Recall@$k$, Precision@$k$, MRR, nDCG@$k$,
Hit@$k$, reported overall and per superclass.\\
\addlinespace
\textbf{Expected values} & Recall@5: \textbf{hybrid RRF 0.939} versus
lexical TF-IDF 0.899. Retrieval runs top-$k$ with $k=5$ against a 67-document
corpus at 768 dimensions under $L_2$.\\
\bottomrule
\end{tabular}
\end{table}

\noindent
Because the corpus is small and hand-seeded, these figures verify that the
retrieval \emph{mechanism} works; they are not evidence of retrieval
quality at realistic corpus scale.

\subsection{Grounding check (local runtime)}

Cell~31 is the acceptance test for the whole chain. It reloads
\code{core/pipeline.py}, calls \code{get\_model()} and asserts three
things: the encoder reports \code{encoder\_is\_trained}, the PTB-XL
diagnostic head is loaded, and the retrieval source is a member of
\code{GROUNDED\_RETRIEVAL\_SOURCES} (\code{hybrid-rrf} or
\code{milvus-sbert}, \emph{not} \code{demo-fallback}). Membership is
tested rather than equality against a literal, so that a change of
retrieval strategy does not silently report correctly-retrieved documents
as ungrounded.

%=============================================================================
\section{Consolidated Results}
%=============================================================================

\begin{table}[!ht]
\centering
\caption{Held-out results, all modules (source: \code{data/genuine\_results.json})}
\label{tab:summary}
\small
\begin{tabular}{@{}p{2.9cm}p{2.7cm}p{2.0cm}p{2.0cm}p{2.0cm}p{1.9cm}@{}}
\toprule
\textbf{Module} & \textbf{Dataset} & \textbf{Accuracy} &
\textbf{Macro-$F_1$} & \textbf{AUC} & \textbf{Trained}\\
\midrule
Flagship ECG & PTB-XL & 0.5778 & 0.4866 & 0.8007 & \yes\\
Beat arrhythmia & MIT-BIH & 0.6065 & 0.5086 & 0.9283 & \yes\\
Ventricular (CV) & VFDB & 0.8672 & 0.7184$^{\dagger}$ & \textbf{0.9187} & \yes\\
Heart sound & CirCor & 0.7604 & 0.5347 & 0.7611$^{\ddagger}$ & \yes\\
Heart disease & UCI Cleveland & 0.7967 & 0.7939 & 0.8718 & \yes\\
Arthritis & NHANES & 0.7622 & 0.6844 & 0.7925 & \yes\\
CardioFusion & MIT-BIH (ECG) & 0.8944 & 0.5401 & 0.8384 & partial\\
LVEF & EchoNet & --- & --- & --- & \no\\
\midrule
RAG retrieval & 67 docs / 38 queries &
\multicolumn{3}{l}{Recall@5 0.939 (hybrid RRF) vs 0.899 (lexical)} & \yes\\
\bottomrule
\multicolumn{6}{@{}l}{\footnotesize $\dagger$ $F_1$(dangerous), the
safety-relevant binary head. $\ddagger$ Present-vs-Absent AUC.}\\
\multicolumn{6}{@{}l}{\footnotesize CardioFusion: only the arrhythmia head
is trained; five task heads remain at initialisation.}
\end{tabular}
\end{table}

\noindent
\textbf{What this table does not say.} No module here exceeds its published
benchmark. The flagship encoder at AUC 0.8007 sits below the
$\approx$0.92 reported for PTB-XL in the literature; the murmur classifier
at macro-$F_1$ 0.5347 is below the CirCor challenge figures; the heart
disease deep model is beaten by logistic regression on its own folds.
Report-generation quality (BLEU, factuality) has never been measured and
no figure for it should be quoted.

%=============================================================================
\section{Reproduction Procedure}
%=============================================================================

\begin{enumerate}[leftmargin=1.5em,itemsep=2pt]
  \item \textbf{Colab, T4 GPU.} Run cell~2 (mounts Drive, detects runtime),
        then cell~4 (dependencies), then cell~5 and confirm in-kernel
        \code{pandas 2.2.3 | wfdb 4.3.1 | numpy 2.0.2}. Pip's own output is
        not sufficient evidence --- a pip upgrade lands underneath a kernel
        that has already imported the old version.
  \item Cell~7 (Kaggle credentials), cell~9 (PTB-XL completeness), cell~11
        (VFDB + CirCor). Datasets already on Drive are not re-downloaded;
        \code{wget -N -c} skips what is current.
  \item Either cell~15 (\code{run\_all}) or the individual module cells
        17--19, 21, 23--26. Every step is guarded: a missing dataset
        records an \code{error} entry and the remaining modules continue.
  \item Cell~37 for VFDB cross-validation (10 fits, $\approx$177~s).
  \item Cell~35 for the ablation studies.
  \item Wait for Drive to finish syncing \code{data/} back down.
  \item \textbf{Switch to the local runtime}, re-run cell~2 \emph{only},
        then cell~20 (pinned client), cells~28--29 (Milvus checks),
        cell~31 (grounding check), cell~33 (RAG evaluation).
\end{enumerate}

\noindent
Prerequisites on the local host: Docker Desktop running with the
\code{milvus-standalone} container started (it does not auto-start after a
reboot), and the Jupyter bridge installed
(\code{jupyter\_http\_over\_ws} plus \code{notebook}).

%=============================================================================
\appendix
\section{Paper-versus-Code Specification Deltas}
\label{app:delta}
%=============================================================================

While assembling this document from source, the following architectural
values in \code{DeepCardio\_CompIntel\_Submit.tex} were found to differ
from the implementation. They are recorded here rather than silently
reconciled; the \emph{code} column is what actually ran to produce the
results in Table~\ref{tab:summary}.

\begin{table}[!ht]
\centering
\footnotesize
\begin{tabular}{@{}p{2.8cm}p{3.7cm}p{3.3cm}p{4.8cm}@{}}
\toprule
\textbf{Item} & \textbf{Paper states} & \textbf{Code implements} &
\textbf{Source}\\
\midrule
Heart-disease BERT-MoE & 4-layer BERT, 6-expert MoE &
\textbf{2 layers, 4 experts} &
\code{heart\_disease\_}\newline\code{pipeline.py:289}\\
\addlinespace
Arthritis BERT-MoE & 4-layer BERT, 6 experts &
\textbf{3 layers}, 6 experts &
\code{arthritis\_pipeline.py:704}\\
\addlinespace
PCG Mel bands & 128 Mel bins & \textbf{64} &
\code{heart\_sound\_loader.py:47}\\
\addlinespace
PCG hop length & hop 256 & \textbf{128} &
\code{heart\_sound\_loader.py:48}\\
\addlinespace
PCG window & 3-second, 50\% overlap &
\textbf{5-second} segments &
\code{heart\_sound\_loader.py:46}\\
\addlinespace
Echo video depth & $64\times112\times112$ &
\textbf{32}$\times112\times112$ &
\code{echonet\_loader.py:27}\\
\addlinespace
Beat-classifier input & ``two leads: Lead~II + modified~V1'' &
single-channel tensor $(N,1,187)$ &
\code{ecg\_arrhythmia\_}\newline\code{pipeline.py:185}\\
\bottomrule
\end{tabular}
\end{table}

\noindent
A separate stale descriptor exists in the code itself:
\code{arthritis\_pipeline.py:1127} still builds the string
\code{"Hybrid Stack: BERT-MoE + GBM + RF + ExtraTrees + SVC $\rightarrow$
LogReg Meta-Learner"}, although the fitted base learners are only
\{\code{gbm}, \code{rf}, \code{et}\}. The same phantom SVC was removed from
the paper and from the dashboard; this occurrence was not part of either
fix.

%=============================================================================
\section{Artefacts Written to \texorpdfstring{\code{data/}}{data/}}
%=============================================================================

% Deliberately NOT a float: as a float it was deferred and dumped at the top
% of the page, i.e. above its own section heading.
\begin{center}
\footnotesize
\begin{tabular}{@{}p{6.6cm}p{8.0cm}@{}}
\toprule
\textbf{File} & \textbf{Written by / consumed by}\\
\midrule
\code{ecg\_ptbxl\_encoder.pt} & \code{train\_ptbxl} $\rightarrow$
\code{core/pipeline.py} (flagship inference)\\
\code{ecg\_ptbxl\_metrics.json} & \code{train\_ptbxl}\\
\code{ecg\_arrhythmia\_cnn\_model.pt},\newline\code{ecg\_arrhythmia\_scaler.pkl} &
\code{train\_ecg\_arrhythmia}\\
\code{vfdb\_detector.pt}, \code{vfdb\_metrics.json} & \code{train\_vfdb}\\
\code{vfdb\_cv\_metrics.json} & \code{cross\_validate\_vfdb}\\
\code{heart\_sound\_classifier.pt},\newline\code{circor\_metrics.json} &
\code{train\_circor}\\
\code{heart\_disease\_bert\_moe\_model.pt},\newline\code{heart\_disease\_bert\_moe\_scaler.pkl},
\newline\code{heart\_disease\_metrics.json} & \code{train\_heartdisease}\\
\code{arthritis\_bert\_moe\_model.pt},\newline\code{arthritis\_bert\_moe\_scaler.pkl} &
\code{train\_arthritis}; the scaler carries \code{feature\_cols}\\
\code{cardiofusion\_weights.pt},\newline\code{cardiofusion\_metrics.json} &
\code{train\_cardiofusion}\\
\code{genuine\_results.json} & \code{run\_all} --- the single source of
truth for every reported metric\\
\bottomrule
\end{tabular}
\end{center}

\vspace{1em}
\begin{notebox}
\textbf{Standing rule.} \code{data/genuine\_results.json} is the only
sanctioned source of performance figures for this system. A number that
cannot be traced to it, or to the raw metrics files it summarises, should
not appear in a paper, a dashboard tile, or a report.
\end{notebox}

\end{document}
